% ==============================================================================
% reports/chapters/02_co_so_ly_thuyet.tex
% CHƯƠNG 2: CƠ SỞ LÝ THUYẾT
% ==============================================================================

\chapter{Cơ Sở Lý Thuyết}
\label{chap:co_so_ly_thuyet}

\section{Bản chất lưu lượng mạng và cơ chế phân tích luồng}

\subsection{Khái niệm luồng mạng}
Theo chuẩn RFC 7011 (IPFIX) \cite{rfc7011}, một luồng mạng được định nghĩa là một tập hợp các gói tin hoặc khung truyền đi qua một điểm quan sát trong một khoảng thời gian nhất định và chia sẻ chung một bộ thuộc tính nhận dạng đặc trưng. Trong thực tế và trong bối cảnh bài tập lớn này, bộ 5 thuộc tính nhận dạng truyền thống là cách nhận dạng flow phổ biến được sử dụng:
\begin{equation}
    \text{Flow 5-Tuple} = \langle \text{Source IP}, \text{Destination IP}, \text{Source Port}, \text{Destination Port}, \text{Protocol} \rangle
\end{equation}

\subsection{Cơ chế sinh đặc trưng của CICFlowMeter}
Công cụ CICFlowMeter \cite{sharafaldin2018generating} thực hiện tổng hợp các gói tin thuộc cùng một 5-Tuple theo cả hai hướng truyền thông (hướng đi từ máy khách đến máy chủ, và hướng về từ máy chủ tới máy khách) để tạo thành một luồng hai chiều. Một flow được coi là kết thúc khi xảy ra một trong các điều kiện sau:
\begin{enumerate}
    \item Xuất hiện cờ điều khiển đóng kết nối TCP (\texttt{FIN} hoặc \texttt{RST}).
    \item Khoảng thời gian không có gói tin mới vượt quá ngưỡng chờ không hoạt động (mặc định 15 đến 60 giây).
    \item Tổng thời lượng của flow vượt quá thời gian tối đa cho phép (Flow Timeout).
\end{enumerate}

Sau khi flow đóng, CICFlowMeter tính toán các chỉ số thống kê tổng hợp mô tả toàn bộ vòng đời của flow, bao gồm thời lượng phiên truyền, thống kê kích thước gói tin (min, max, mean, std), thời gian trễ giữa các gói tin liên tiếp (Flow IAT, Fwd IAT, Bwd IAT), các cờ điều khiển TCP (SYN, ACK, PSH, URG, FIN) và kích thước cửa sổ nhận ban đầu (\texttt{Init\_Win\_bytes}).

\subsection{Phân loại luồng đã hoàn tất và NIDS thời gian thực}
Một điểm cốt lõi cần làm rõ về mặt lý thuyết: vì toàn bộ các đặc trưng thống kê nói trên chỉ được tính toán đầy đủ \textbf{sau khi flow đã kết thúc hoàn toàn}, các mô hình học máy được xây dựng trên dữ liệu này hoạt động ở cơ chế \textbf{Phân loại luồng mạng đã hoàn tất}.

Hệ thống này \textbf{không phải} là hệ thống phát hiện xâm nhập thời gian thực ở mức gói tin. Mô hình không can thiệp ngăn chặn tấn công ngay ở những gói tin bắt tay đầu tiên và không đo lường độ trễ phát hiện theo mili-giây.

\section{Cơ chế tấn công Brute Force trên giao thức FTP và SSH}

\subsection{Tấn công Brute Force đối với FTP (Cổng 21)}
FTP là giao thức truyền tệp hoạt động trên kênh văn bản rõ (plain-text) ở tầng ứng dụng:
\begin{itemize}
    \item Quá trình xác thực diễn ra thông qua các cặp lệnh \texttt{USER} và \texttt{PASS}.
    \item Máy chủ phản hồi bằng mã trạng thái \texttt{230 User logged in} (thành công) hoặc \texttt{530 Login incorrect} (thất bại).
    \item Khi công cụ Patator (\texttt{ftp\_patator}) thực hiện tấn công dò quét mật khẩu, nó thiết lập liên tục các kết nối TCP mới, gửi lệnh thử mật khẩu và đóng kết nối ngay khi nhận mã 530.
    \item \textbf{Đặc trưng thống kê}: Vì là văn bản rõ, các gói tin yêu cầu và phản hồi có kích thước ngắn và tương đối cố định. Lưu lượng hướng đi (Fwd) thường có độ lệch chuẩn kích thước gói tin nhỏ.
\end{itemize}

\subsection{Tấn công Brute Force đối với SSH (Cổng 22)}
SSH là giao thức truyền thông an toàn cung cấp kênh xác thực và điều khiển được bảo vệ bằng mật mã:
\begin{itemize}
    \item Trước khi quá trình gửi thông tin xác thực diễn ra, client và server trải qua pha bắt tay mật mã: trao đổi phiên bản giao thức, thương lượng thuật toán mã hóa đối xứng, trao đổi khóa và xác thực khóa công khai của host.
    \item Quá trình xác thực mật khẩu diễn ra hoàn toàn bên trong kênh đã được bảo vệ.
    \item Đặc trưng thống kê: Pha trao đổi khóa mật mã làm phát sinh các gói tin có kích thước lớn hơn đáng kể so với FTP. Cấu trúc bắt tay khiến các chỉ số về kích thước gói tin trung bình (\texttt{Packet Length Mean}) và thời gian trễ giữa các gói tin (\texttt{Flow IAT}) của SSH có phân phối thống kê khác biệt so với FTP.
\end{itemize}

\section{Các thuật toán học máy dạng cây}

\subsection{Cây quyết định (Decision Tree)}
Decision Tree \cite{quinlan1986induction} phân chia không gian đặc trưng thành các vùng siêu hình hộp rời rạc thông qua một tập hợp các quy tắc rẽ nhánh dạng bậc thang ($x_j \le \theta$).

Tại mỗi nút phân chia, thuật toán tìm kiếm đặc trưng $j$ và ngưỡng $\theta$ tối ưu nhằm tối đa hóa mức giảm độ vẩn đục. Chỉ số vẩn đục phổ biến là \textbf{Gini Impurity}:
\begin{equation}
    I_G(p) = 1 - \sum_{k=0}^{1} p_k^2
\end{equation}
trong đó $p_k$ là tỷ lệ mẫu thuộc lớp $k$ tại nút đang xét. Ưu điểm nổi bật của Decision Tree là tính minh bạch và khả năng trực quan hóa cấu trúc rẽ nhánh, giúp giải thích logic phân loại. Tuy nhiên, Decision Tree đơn lẻ có phương sai cao và dễ bị quá khớp.

\subsection{Rừng ngẫu nhiên (Random Forest)}
Random Forest \cite{breiman2001random} là phương pháp học kết hợp dựa trên kỹ thuật đóng bao. Mô hình xây dựng một tập hợp gồm $B$ cây quyết định độc lập:
\begin{equation}
    \hat{y} = \text{mode} \left\{ T_1(\mathbf{x}), T_2(\mathbf{x}), \dots, T_B(\mathbf{x}) \right\}
\end{equation}

Hai cơ chế ngẫu nhiên hóa cốt lõi của Random Forest giúp giảm phương sai mà không làm tăng độ chệch:
\begin{enumerate}
    \item \textbf{Bootstrap Sampling}: Mỗi cây được huấn luyện trên một tập dữ liệu con được rút mẫu có hoàn lại từ tập huấn luyện gốc.
    \item \textbf{Random Feature Subspace}: Tại mỗi nút rẽ nhánh, thuật toán chỉ chọn ngẫu nhiên một tập con gồm $m \approx \sqrt{D}$ đặc trưng để chọn ra điểm phân chia tốt nhất, giúp giải tương quan giữa các cây.
\end{enumerate}

\subsection{XGBoost}
XGBoost \cite{chen2016xgboost} là thuật toán nâng cao thuộc họ Gradient Boosting. Boosting xây dựng các cây quyết định tuần tự, trong đó mỗi cây mới $f_t(\mathbf{x})$ được tối ưu để bù đắp sai số của tổ hợp các cây trước đó:
\begin{equation}
    \hat{y}_i^{(t)} = \hat{y}_i^{(t-1)} + f_t(\mathbf{x}_i)
\end{equation}

Hàm mục tiêu được tối thiểu hóa tại bước $t$ bao gồm hàm mất mát khả vi và số hạng chính quy hóa để kiểm soát độ phức tạp của cây:
\begin{equation}
    \mathcal{L}^{(t)} = \sum_{i=1}^{n} l\left(y_i, \hat{y}_i^{(t-1)} + f_t(\mathbf{x}_i)\right) + \Omega(f_t)
\end{equation}
với số hạng phạt:
\begin{equation}
    \Omega(f) = \gamma T + \frac{1}{2} \lambda \sum_{j=1}^{T} w_j^2
\end{equation}
trong đó $T$ là số lượng lá, $w_j$ là trọng số tại lá $j$, $\gamma$ là ngưỡng phạt tối thiểu để tạo lá mới và $\lambda$ là hệ số chính quy hóa L2.

\subsection{Khả năng giải thích mô hình với SHAP}
Để giải thích đóng góp của từng đặc trưng vào dự đoán của các mô hình dạng cây, phương pháp SHAP \cite{lundberg2017unified} dựa trên lý thuyết trò chơi hợp tác Shapley được sử dụng. Giá trị SHAP gán cho mỗi đặc trưng đại diện cho mức độ thay đổi kỳ vọng dự đoán khi đặc trưng đó xuất hiện so với khi vắng mặt:
\begin{equation}
    \phi_i(x) = \sum_{S \subseteq F \setminus \{i\}} \frac{|S|!(|F| - |S| - 1)!}{|F|!} \left[ f_x(S \cup \{i\}) - f_x(S) \right]
\end{equation}
Phương pháp TreeExplainer cho phép tính toán chính xác giá trị SHAP trên các mô hình dạng cây với độ phức tạp đa thức, giúp xác định các đặc trưng có ảnh hưởng lớn nhất đến quyết định phân loại.

\section{Đánh giá mô hình trên dữ liệu mất cân bằng}

\subsection{Hệ thống chỉ số từ ma trận nhầm lẫn}
Trong bài toán phát hiện xâm nhập mạng, lớp tấn công (\Attack, nhãn 1) thường chiếm tỷ lệ nhỏ so với lưu lượng bình thường (\Normal, nhãn 0). Ma trận nhầm lẫn gồm 4 đại lượng: True Positive ($TP$), False Positive ($FP$), False Negative ($FN$) và True Negative ($TN$).

Các chỉ số phái sinh chuẩn:
\begin{equation}
    \text{Precision} = \frac{TP}{TP + FP}, \quad \text{Recall} = \frac{TP}{TP + FN}
\end{equation}
\begin{equation}
    \text{False Positive Rate (FPR)} = \frac{FP}{FP + TN}
\end{equation}
\begin{equation}
    \text{F1-Score} = 2 \cdot \frac{\text{Precision} \cdot \text{Recall}}{\text{Precision} + \text{Recall}} = \frac{2TP}{2TP + FP + FN}
\end{equation}

\subsection{Nghịch lý độ chính xác trong NIDS}
Chỉ số Accuracy tổng thể:
\begin{equation}
    \text{Accuracy} = \frac{TP + TN}{TP + TN + FP + FN}
\end{equation}
có thể gây hiểu nhầm nghiêm trọng nếu được sử dụng độc lập trên dữ liệu mất cân bằng. Giả sử tập kiểm thử có $97.4\%$ mẫu bình thường và $2.6\%$ mẫu tấn công (như tập Test CICIDS2017 Tuesday). Một bộ phân loại tầm thường luôn dự đoán mọi flow là bình thường ($TP=0, FP=0$) sẽ đạt ngay Accuracy = 97.4\% mặc dù không phát hiện được bất kỳ cuộc tấn công nào (Recall = 0.0\%). Do đó, F1-Score, đường cong Precision-Recall và Average Precision là các chỉ số đánh giá cần thiết.

\subsection{Đường cong Precision-Recall và Average Precision}
Như Davis và Goadrich \cite{davis2006relationship} đã chỉ ra, trên các tập dữ liệu mất cân bằng lớp, đường cong ROC có thể tạo ra cảm nhận quá lạc quan do đại lượng $TN$ trong mẫu số của $FPR$ rất lớn. Đường cong Precision-Recall biểu diễn sự đánh đổi giữa Precision và Recall khi thay đổi ngưỡng phân loại $\tau \in [0, 1]$.

Theo đặc tả của scikit-learn \cite{scikit_learn_ap}, chỉ số Average Precision được tính bằng giá trị trung bình có trọng số của Precision theo mức tăng Recall giữa các ngưỡng:
\begin{equation}
    \text{AP} = \sum_{n} (R_n - R_{n-1}) P_n
\end{equation}
trong đó $P_n$ và $R_n$ là Precision và Recall tại ngưỡng thứ $n$. AP đánh giá chất lượng xếp hạng các mẫu dương trên dải ngưỡng phân loại; AP không đồng nhất hoàn toàn với diện tích tích phân hình thang thô của PR-AUC và không đo lường độ chuẩn xác xác suất.
